\documentclass{beamer}
\usepackage{graphicx}
\usepackage{xcolor}
\usepackage{tikz}
\usepackage{amsmath}

\definecolor{purpleblue}{RGB}{80, 80, 180}
\definecolor{novelPurple}{RGB}{98, 54, 255}
\definecolor{softPurple}{RGB}{180, 170, 255}
\definecolor{softGray}{RGB}{235, 235, 245}
\definecolor{lightText}{RGB}{90, 90, 110}

\usetheme{Madrid}
\usecolortheme{dolphin}
\graphicspath{{static/course_materials/}{static/hard/}{./}}

\setbeamerfont{title}{size=\Large,series=\bfseries}
\setbeamerfont{frametitle}{size=\LARGE,series=\bfseries}
\setbeamercolor{frametitle}{fg=white, bg=novelPurple}
\setbeamercolor{title}{fg=white, bg=purpleblue}

\title[SAT Hard Question 6]{\textbf{SAT Hard Question\\ Set 6}}
\author{\textbf{Jack Zeng}}
\institute{\textcolor{white}{\textbf{Novel Prep}}}
\date{\today}

\begin{document}

\begin{frame}
    \titlepage
    \vfill
    \begin{flushright}
        \textcolor{purpleblue}{\Huge \textbf{Novel Prep}}
    \end{flushright}
\end{frame}

\begin{frame}
    \begin{tikzpicture}[remember picture, overlay]
        \shade[top color=softPurple, bottom color=softGray]
            (current page.north west) rectangle (current page.south east);
        \node[align=center] at (current page.center) {
            {\Huge\bfseries\textcolor{purpleblue}{Novel Prep}} \\[0.5cm]
            {\Large\itshape\textcolor{lightText}{Phase 2 · Hard Question Practice}}
        };
    \end{tikzpicture}
\end{frame}

\begin{frame}{Overview}
    \tableofcontents
\end{frame}

\begin{frame}
\section{Hard Question Set 6}
    \begin{tikzpicture}[remember picture, overlay]
        \shade[top color=softPurple, bottom color=softGray]
            (current page.north west) rectangle (current page.south east);
        \node[align=center] at (current page.center) {
            {\LARGE\bfseries\textcolor{purpleblue}{Hard Question Set 6}} \\[0.5cm]
            {\Large\itshape\textcolor{lightText}{15 SAT-style challenge problems}}
        };
    \end{tikzpicture}
\end{frame}

% --- Question 1 ---
\begin{frame}{Question 1}
\small
Given the expression
\[
\sqrt[5]{70n \left(\sqrt[7]{70n}\right)^2}
\]
for what value of \(x\) is the given expression equivalent to \((70n)^{30x}\), where \(n > 1\)?
\end{frame}
\begin{frame}{Answer 1}
\small
\textbf{Final Answer:} $\boxed{\dfrac{3}{350}}$
\end{frame}

% --- Question 2 ---
\begin{frame}{Question 2}
\small
The expression \(4x^2 + bx - 45\), where \(b\) is a constant, can be rewritten as
\[
(hx + k)(x + j)
\]
where \(h, k,\) and \(j\) are integer constants.

Which of the following must be an integer?
\vspace{0.35cm}
\begin{enumerate}
    \item[A.] \(\dfrac{b}{h}\)
    \item[B.] \(\dfrac{b}{k}\)
    \item[C.] \(\dfrac{45}{h}\)
    \item[D.] \(\dfrac{45}{k}\)
\end{enumerate}
\end{frame}
\begin{frame}{Answer 2}
\small
\textbf{Correct Answer:} D
\end{frame}

% --- Question 3 ---
\begin{frame}{Question 3}
\small
The expression
\[
6\sqrt[5]{3x^{45}} \cdot \sqrt[8]{8x}
\]
is equivalent to \(ax^b\), where \(a\) and \(b\) are positive constants and \(x > 1\).

What is the value of \(b\)?
\end{frame}
\begin{frame}{Answer 3}
\small
\textbf{Final Answer:} $\boxed{\dfrac{73}{8}}$
\end{frame}

% --- Question 4 ---
\begin{frame}{Question 4}
\small
In the given equation, \(c\) is a positive constant. Solve for \(x\):
\[
\frac{x^2}{\sqrt{x^2 - c^2}} = \frac{c^2}{\sqrt{x^2 - c^2}} + 39
\]

Which of the following is one of the solutions to the given equation?
\vspace{0.35cm}
\begin{enumerate}
    \item[A.] \(-c\)
    \item[B.] \(-c^2 - 39^2\)
    \item[C.] \(-\sqrt{39^2 - c^2}\)
    \item[D.] \(\sqrt{c^2 + 39^2}\)
\end{enumerate}
\end{frame}
\begin{frame}{Answer 4}
\small
\textbf{Correct Answer:} D
\end{frame}

% --- Question 5 ---
\begin{frame}{Question 5}
\small
The quadratic function is given as
\[
f(x) = ax^2 + 4x + c
\]
where \(a\) and \(c\) are constants. The graph of \(y = f(x)\) is a parabola that opens upward and has a vertex at \((h, k)\), where \(h\) and \(k\) are constants.

If \(k < 0\) and \(f(-9) = f(3)\), which of the following must be true?

\begin{itemize}
\item[I] \(c < 0\)
\item[II] \(a \geq 1\)
\end{itemize}
\vspace{0.35cm}
\begin{enumerate}
    \item[A.] I only
    \item[B.] II only
    \item[C.] I and II
    \item[D.] Neither I nor II
\end{enumerate}
\end{frame}
\begin{frame}{Answer 5}
\small
\textbf{Correct Answer:} D
\end{frame}

% --- Question 6 ---
\begin{frame}{Question 6}
\small
The functions \(f\) and \(g\) are defined by the given equations, where \(x \geq 0\):

\begin{itemize}
\item[I] \(f(x) = 33(0.4)^{x+3}\)
\item[II] \(g(x) = 33(0.16)(0.4)^{x-2}\)
\end{itemize}

Which of the following equations displays, as a constant or coefficient, the maximum value of the function it defines, where \(x \geq 0\)?
\vspace{0.35cm}
\begin{enumerate}
    \item[A.] I only
    \item[B.] II only
    \item[C.] I and II
    \item[D.] Neither I nor II
\end{enumerate}
\end{frame}
\begin{frame}{Answer 6}
\small
\textbf{Correct Answer:} B
\end{frame}

% --- Question 7 ---
\begin{frame}{Question 7}
\small
Each of the following frequency tables represents a data set.

Which of these frequency tables represents the data set with the \textbf{smallest standard deviation}?

\[
\begin{array}{|c|c|c|c|c|}
\hline
\text{Value} & \text{Frequency (A)} & \text{Frequency (B)} & \text{Frequency (C)} & \text{Frequency (D)} \\
\hline
5 & 0 & 11 & 15 & 18 \\
6 & 10 & 17 & 15 & 14 \\
7 & 55 & 19 & 15 & 11 \\
8 & 10 & 17 & 15 & 14 \\
9 & 0 & 11 & 15 & 18 \\
\hline
\end{array}
\]
\vspace{0.35cm}
\begin{enumerate}
    \item[A.] A
    \item[B.] B
    \item[C.] C
    \item[D.] D
\end{enumerate}
\end{frame}
\begin{frame}{Answer 7}
\small
\textbf{Correct Answer:} A
\end{frame}

% --- Question 8 ---
\begin{frame}{Question 8}
\small
In \(\triangle ABC\), \(AB = 4\) and \(BC = 5\). If \(\angle ABC < 90^\circ\), which of the following could be the area of \(\triangle ABC\)?
\vspace{0.35cm}
\begin{enumerate}
    \item[A.] 5
    \item[B.] 10
    \item[C.] 15
    \item[D.] 20
\end{enumerate}
\end{frame}
\begin{frame}{Answer 8}
\small
\textbf{Correct Answer:} A
\end{frame}

% --- Question 9 ---
\begin{frame}{Question 9}
\small
The table summarizes the distribution of heights, in inches (\("in"\)), for 300 participants in a study.

If a participant that is at most 71 inches tall is chosen at random, what is the probability that the participant was selected from Group 2?

\[
\begin{array}{|c|c|c|c|c|}
\hline
 & \text{48--59 in.} & \text{60--71 in.} & \text{72--83 in.} & \text{Total} \\
\hline
\text{Group 1} & 20 & 60 & 20 & 100 \\
\text{Group 2} & 10 & 75 & 15 & 100 \\
\text{Group 3} & 30 & 45 & 25 & 100 \\
\hline
\text{Total} & 60 & 180 & 60 & 300 \\
\hline
\end{array}
\]
\vspace{0.35cm}
\begin{enumerate}
    \item[A.] \(\dfrac{1}{4}\)
    \item[B.] \(\dfrac{17}{48}\)
    \item[C.] \(\dfrac{17}{60}\)
    \item[D.] \(\dfrac{17}{20}\)
\end{enumerate}
\end{frame}
\begin{frame}{Answer 9}
\small
\textbf{Correct Answer:} B
\end{frame}

% --- Question 10 ---
\begin{frame}{Question 10}
\small
The function \(f\) is defined as \(f(x) = x^2 + ax + b\), where \(a\) and \(b\) are positive integers.

If the minimum value of \(f(x)\) is 7, which is a possible value of \(b\)?
\vspace{0.35cm}
\begin{enumerate}
    \item[A.] 6
    \item[B.] 7
    \item[C.] 8
    \item[D.] 9
\end{enumerate}
\end{frame}
\begin{frame}{Answer 10}
\small
\textbf{Correct Answer:} C
\end{frame}

% --- Question 11 ---
\begin{frame}{Question 11}
\small
The average annual energy cost for a certain home is \$4{,}334. The homeowner plans to spend \$25{,}000 to install a geothermal heating system. The homeowner estimates that the average annual energy cost will then be \$2{,}712.

Which of the following inequalities can be solved to find \(t\), the number of years after installation at which the total amount of energy cost savings will exceed the installation cost?
\vspace{0.35cm}
\begin{enumerate}
    \item[A.] \(25{,}000 > (4{,}334 - 2{,}712)t\)
    \item[B.] \(25{,}000 < (4{,}334 - 2{,}712)t\)
    \item[C.] \(25{,}000 - 4{,}334 - 2{,}712t\)
    \item[D.] \(25{,}000 > \dfrac{4{,}332}{2{,}712}t\)
\end{enumerate}
\end{frame}
\begin{frame}{Answer 11}
\small
\textbf{Correct Answer:} B
\end{frame}

% --- Question 12 ---
\begin{frame}{Question 12}
\small
For the equation, \(b\) and \(c\) are constants. For any real number \(t\), the point
\[
\left(\frac{-2t + 4}{3},\, t\right)
\]
lies on the graph of the equation \(bx + 3y = c\).

What is the value of \(c\)?
\vspace{0.35cm}
\begin{enumerate}
    \item[A.] -2
    \item[B.] 4
    \item[C.] 6
    \item[D.] 12
\end{enumerate}
\end{frame}
\begin{frame}{Answer 12}
\small
\textbf{Correct Answer:} C
\end{frame}

% --- Question 13 ---
\begin{frame}{Question 13}
\small
A circle in the \(xy\)-plane has equation
\[
(x - 2.5)^2 + (y - 3.5)^2 = 225.
\]

Line \(p\) is tangent to the circle at point \((-6.5, C)\), where \(C\) is a constant.

Which of the following could be the slope of line \(p\)?

\begin{itemize}
\item[I] \(-\dfrac{3}{4}\)
\item[II] \(\dfrac{3}{4}\)
\item[III] \(\dfrac{4}{3}\)
\end{itemize}
\vspace{0.35cm}
\begin{enumerate}
    \item[A.] I only
    \item[B.] I and II only
    \item[C.] I and III only
    \item[D.] I, II, and III
\end{enumerate}
\end{frame}
\begin{frame}{Answer 13}
\small
\textbf{Correct Answer:} B
\end{frame}

% --- Question 14 ---
\begin{frame}{Question 14}
\small
The function \(f\) is defined by \(f(x) = a^x - b\), where \(a\) and \(b\) are constants.

In the \(xy\)-plane, the graph of \(y = f(x)\) passes through the points \((c, 7)\) and \((2c, 117)\), where \(c\) is a constant.

What is the value of \(b\)?
\end{frame}
\begin{frame}{Answer 14}
\small
\textbf{Final Answer:} $\boxed{4}$
\end{frame}

% --- Question 15 ---
\begin{frame}{Question 15}
\small
The given equation, where \(t\) is a positive constant, defines a circle in the \(xy\)-plane.

The radius of this circle is \(\sqrt{41}\).

What is the value of \(t\)?
\[
2x^2 + 2y^2 + tx + ty - 78 = 0
\]
\end{frame}
\begin{frame}{Answer 15}
\small
\textbf{Final Answer:} $\boxed{4}$
\end{frame}

\begin{frame}{}
    \begin{tikzpicture}[remember picture, overlay]
        \shade[top color=novelPurple!40, bottom color=white]
            (current page.north west) rectangle (current page.south east);
        \fill[white, opacity=0.15] (current page.center) circle (5cm);
        \foreach \x in {1,2,3,4,5} {
            \fill[novelPurple, opacity=0.12] (rand*10-5, rand*7-3.5) circle (0.1+\x*0.05);
        }
    \end{tikzpicture}
    \centering
    \vspace{5em}
    {\Huge \textbf{\textcolor{novelPurple}{Thank You!}}} \\[1em]
    {\large \textcolor{gray}{We appreciate your time and focus.}} \\[3em]
    {\LARGE \textbf{\textcolor{novelPurple}{\textsf{Novel Prep}}}}
    \vfill
\end{frame}

\end{document}
